\documentclass[a4paper]{article} 
\addtolength{\hoffset}{-2.25cm}
\addtolength{\textwidth}{4.5cm}
\addtolength{\voffset}{-3.25cm}
\addtolength{\textheight}{5cm}
\setlength{\parskip}{0pt}
\setlength{\parindent}{0in}

%----------------------------------------------------------------------------------------
%	PACKAGES AND OTHER DOCUMENT CONFIGURATIONS
%----------------------------------------------------------------------------------------
\usepackage{graphicx}
\usepackage{listings}
\graphicspath{{img/}}
\usepackage[section]{placeins}
\usepackage{blindtext} % Package to generate dummy text
\usepackage{charter} % Use the Charter font
\usepackage[utf8]{inputenc} % Use UTF-8 encoding
\usepackage{microtype} % Slightly tweak font spacing for aesthetics
\usepackage[portuguese]{babel} % Language hyphenation and typographical rules
\usepackage{amsthm, amsmath, amssymb} % Mathematical typesetting
\usepackage{float} % Improved interface for floating objects
\usepackage[final, colorlinks = true, 
            linkcolor = black, 
            citecolor = black]{hyperref} % For hyperlinks in the PDF
\usepackage{graphicx, multicol} % Enhanced support for graphics
\usepackage{xcolor} % Driver-independent color extensions
\usepackage{marvosym, wasysym} % More symbols
\usepackage{rotating} % Rotation tools
\usepackage{censor} % Facilities for controlling restricted text
\usepackage{listings, style/lstlisting} % Environment for non-formatted code, !uses style file!
\usepackage{pseudocode} % Environment for specifying algorithms in a natural way
\usepackage{style/avm} % Environment for f-structures, !uses style file!
\usepackage{booktabs} % Enhances quality of tables
\usepackage{enumitem}
\usepackage{float}
\usepackage[backend=biber,style=numeric,
            sorting=nyt]{biblatex} % Complete reimplementation of bibliographic facilities
\addbibresource{ecl.bib}
\usepackage{csquotes} % Context sensitive quotation facilities
\usepackage[yyyymmdd]{datetime} % Uses YEAR-MONTH-DAY format for dates
\renewcommand{\dateseparator}{-} % Sets dateseparator to '-'
\usepackage{fancyhdr} % Headers and footers
\pagestyle{fancy} % All pages have headers and footers
\fancyhead{}\renewcommand{\headrulewidth}{0pt} % Blank out the default header
\fancyfoot[L]{} % Custom footer text
\fancyfoot[C]{} % Custom footer text
\fancyfoot[R]{\thepage} % Custom footer text
\newcommand{\note}[1]{\marginpar{\scriptsize \textcolor{red}{#1}}} % Enables comments in red on margin

%----------------------------------------------------------------------------------------

\begin{document}

% TITLE SECTION

\begin{figure}[h!]
\centering
\includegraphics[width=75mm]{logoDIREC.png}
\end{figure}

\fancyhead[C]{}
\hrule \medskip % Upper rule
\begin{minipage}{0.295\textwidth} 
\raggedright
\footnotesize

Artur Barichello \hfill\\   
Lucas Zacchi \hfill\\   
Monique Bertan \hfill\\

\end{minipage}
\begin{minipage}{0.4\textwidth} 
\centering 
\large 
Relatório 1\\ 
\normalsize 
Sistemas Digitais\\
\end{minipage}
\begin{minipage}{0.295\textwidth} 
\raggedleft
5 de abril de 2018\hfill\\
\end{minipage}
\medskip\hrule 
\bigskip

% CONTENTS
%------------------------------------------------
\section{\textbf{Descrição do problema}}
%------------------------------------------------

Descrever uma ULA para números negativos e positivos de 16 bits com os seguintes componentes:
\begin{itemize}
\item Dois registradores.
\item Uma unidade para cálculo de complemento de 2.
\item Um seletor para utilizar ou pular o uso do complemento de dois e outro seletor anterior ao registrador A.
\item Um registrador de 1 bit para detecção de overflow.
\item Um somador com detecção de overflow.
\end{itemize}

\smallskip

%------------------------------------------------
\section{\textbf{Descrição da elaboração do trabalho}}
%-------------------------

O projeto foi feito com uma arquitetura topo abrigando 6 componentes, 7 instanciações destes componentes e 8 sinais internos entre os blocos. 
Componentes como registradores foram feitos da maneira genérica enquanto os seletores foram feitos do modo estrutural.
\bigskip

\textbf{Mux estrutural}
\begin{lstlisting}[language=VHDL]
LIBRARY ieee;
USE ieee.std_logic_1164.all;
USE ieee.numeric_std.all;

ENTITY mux2x1 IS
GENERIC (n: INTEGER := 16);
PORT (sel  : IN STD_LOGIC;
      w, x : IN SIGNED(n-1 DOWNTO 0);
	  m    : OUT SIGNED(n-1 DOWNTO 0));
END mux2x1;

ARCHITECTURE bhv OF mux2x1 IS
BEGIN
    m <= w WHEN sel = '0' ELSE x;
end bhv;
\end{lstlisting}
\bigskip

O registrador A é um registrador genérico normal enquanto o registrador B possui duas saídas de dados, uma que segue para o complemento de dois enquanto outra saída pula este componente.
\bigskip

\textbf{Registrador A genéricos de 16 bits:}
\begin{lstlisting}[language=VHDL]
LIBRARY ieee;
USE ieee.std_logic_1164.all;
USE ieee.numeric_std.all;

ENTITY register_a IS
GENERIC (n: INTEGER := 16);
PORT (clk : IN STD_LOGIC;
      d   : IN SIGNED(n-1 DOWNTO 0);
      q   : OUT SIGNED(n-1 DOWNTO 0));
END register_a;

ARCHITECTURE bhv OF register_a IS
BEGIN
    PROCESS(clk)
    BEGIN
        IF (RISING_EDGE(clk)) THEN
			q <= d;
        END IF;
    END PROCESS;
END bhv;
\end{lstlisting}

\pagebreak
\textbf{Registrador B com saída para o seletor (números positivos)}
\begin{lstlisting}[language=VHDL]
LIBRARY ieee;
USE ieee.std_logic_1164.all;
USE ieee.numeric_std.all;

ENTITY register_b IS
GENERIC (n: INTEGER := 16);
PORT (clk     : IN STD_LOGIC;
      d       : IN SIGNED(n-1 DOWNTO 0);
      q       : OUT SIGNED(n-1 DOWNTO 0);
		sel_out : OUT SIGNED(n-1 DOWNTO 0));
END register_b;

ARCHITECTURE bhv OF register_b IS
BEGIN
    PROCESS(clk)
    BEGIN
        IF (RISING_EDGE(clk)) THEN
				q <= d;
        END IF;
    END PROCESS;
END bhv;
\end{lstlisting}
\smallskip

O módulo do complemento de dois utiliza o operator NOT do quartus que inverte todos os bits da entrada e passa para um sinal interno temporário. Este sinal é incrementado em um e depois é atribuido à saída.
\bigskip

\textbf{Módulo complemento de dois}
\begin{lstlisting}[language=VHDL]
LIBRARY ieee;
USE ieee.std_logic_1164.all;
USE ieee.numeric_std.all;

ENTITY comp2 IS
GENERIC (n: INTEGER := 16);
PORT (a_in  : IN SIGNED(n-1 DOWNTO 0);
		s_out : OUT SIGNED(N-1 DOWNTO 0));
END comp2;

ARCHITECTURE bhv OF comp2 IS
SIGNAL tmp : SIGNED(n-1 DOWNTO 0);

BEGIN
    tmp <= NOT a_in;
	 s_out <= SIGNED(tmp + 1);
end bhv;
\end{lstlisting}
\smallskip

O adder de 16 bits com detecção de overflow foi projetado da seguinte maneira:
É criado um sinal 'sum' com 1 bit a mais que os operandos (no caso 17), onde é feita a soma dos dois números de 16 bits. Se este novo bit à esquerda for 1 foi detectado que houve overflow e o sinal OUT 'overflow' é setado para 1.
\bigskip

\textbf{Adder com detecção de overflow}
\begin{lstlisting}[language=VHDL]
LIBRARY ieee;
USE ieee.std_logic_1164.all;
USE ieee.numeric_std.all;
USE ieee.std_logic_unsigned.all;
USE ieee.std_logic_signed.all;

ENTITY adder IS
PORT (a, b     : IN  SIGNED(15 DOWNTO 0);
      overflow : OUT STD_LOGIC;
	  s        : OUT SIGNED(15 DOWNTO 0));
END adder;

ARCHITECTURE bhv OF adder IS
SIGNAL sum: SIGNED(16 DOWNTO 0);

BEGIN
	PROCESS(sum)
	BEGIN
		sum <= a + b;
		IF sum(16) = '1' THEN
			overflow <= '1';
		ELSE
			overflow <= '0';
			s <= sum;
		END IF;
	END PROCESS;
END bhv;
\end{lstlisting}
\smallskip

\bigskip

%------------------------------------------------
\section{\textbf{Testes}}
%------------------------------------------------

\smallskip

%------------------------------------------------
\section{\textbf{Resultados de atrasos e utilização da placa de cada módulo geral}}
%------------------------------------------------

\smallskip

%------------------------------------------------
\section{\textbf{Analise dos resultados}}
%------------------------------------------------

\smallskip

%------------------------------------------------
\section{\textbf{Arquivos VHDL e Estímulos}}
%------------------------------------------------

% \begin{figure}[H]
% \centering
% \includegraphics[width=90mm]{}
% \caption{Arquivo VHDL \label{Figura 1}}
% \end{figure}
% \smallskip

% END
\bigskip

\end{document}
